%% A Finite Characterization of Plugin-System Conformance for Dynamic Composition
%% ACM acmart format (sigplan). Compile: tectonic main.tex  or  pdflatex+bibtex-free.
\documentclass[sigplan,screen,review,nonacm]{acmart}

\usepackage{booktabs}

%% ---- macros ------------------------------------------------------------
\newcommand{\Spec}{\ensuremath{\mathbb S}}
\newcommand{\Impl}{\ensuremath{\mathbb I}}
\newcommand{\conf}{\mathrel{\mathbf{conf}}}
\newcommand{\CF}{\ensuremath{\mathfrak I_{\mathrm{CF}}}}
\newcommand{\Tstar}{\ensuremath{T^{\star}}}
\newcommand{\occ}[1]{\ensuremath{\Sigma_{#1}}}
\newcommand{\kap}[1]{\ensuremath{\kappa_{#1}}}
\newcommand{\Mfin}{\ensuremath{\mathcal M_{\mathrm{fin}}}}
\newcommand{\target}{\mathrm{target}}
\newcommand{\relied}{\mathrm{relied}}
\newcommand{\Reach}{\mathrm{Reach}}
\newcommand{\Out}{\mathrm{Out}}
\newcommand{\Inactive}{\mathsf{Inactive}}
\newcommand{\Reloading}{\mathsf{Reloading}}
\newcommand{\Active}{\mathsf{Active}}
\newcommand{\Unloading}{\mathsf{Unloading}}
\newcommand{\SZC}{[SZC26]}

\begin{document}

\title{A Finite Characterization of Plugin-System Conformance for Dynamic Composition}

\author{The Cordis Conformance Project}
\affiliation{%
  \institution{Companion artifact to \emph{A Programming Paradigm for Spatiotemporal Composability}}
  \country{}}

\renewcommand{\shortauthors}{Cordis Conformance}

\begin{abstract}
The calculus of dynamic composition of Shi, Zhang, and Cui \SZC{} carries a proved
metatheory---revertible effects, reactive coeffects, a guarded component lifecycle,
and theorems from recovery exactness to confluence---and an implementation, Cordis,
related to the calculus by a correspondence table and reference algorithms, but by
no proof. We supply the missing verification as a theorem about testing. We
formalize the plugin-system interface induced by the paper's own correspondence
table as an instrumented labeled transition system with a fixed observation
vocabulary, and prove first that \emph{no} finite black-box suite characterizes
conformance over unrestricted implementations. We then define the auditable class
of \emph{clause-factored} implementations, whose rule decisions factor through a
finite alphabet of \emph{atomic signatures}---the occurrence-level decomposition of
each rule's window, with unbounded registries entering only through per-key and
per-witness occurrences---and construct, for every reachable signature, a canonical
experiment. The resulting finite suite \Tstar{} characterizes conformance on this
class:
$\Impl \in \CF \Rightarrow (\Impl \conf \Spec \iff \Impl \models \Tstar)$.
The proof is by hand, in the style of the source paper: a factorization lemma
showing the calculus is itself clause-factored; transport and observation-reflection
lemmas (the latter provably \emph{false} without transition certificates: we exhibit
a rule aspect that is unobservable in the event-and-report vocabulary on every
script); a shortest-counterexample safety argument; and a termination argument by
ranking that, en route, \emph{strengthens the source paper's progress theorem by
deleting its acyclicity hypothesis}. The signature tables are enumerated in full
(92 occurrences, 86 reachable, 6 proved unrealizable at the interface). An
exhaustive-schedule test harness grounds the theory: applying the concrete
32-obligation suite to the current Cordis TypeScript implementation verifies one
build on every obligation under complete enumeration of each scenario's schedule
tree, convicts the published build on exactly four (one defect previously unknown),
and finds zero divergence attributable to a two-node distributed deployment. The
construction also returns three corrections to the source paper itself.
\end{abstract}

\maketitle

%% =======================================================================
\section{Introduction}\label{sec:intro}

Dynamic composition frameworks promise that components can be loaded, unloaded,
and reconfigured at runtime while the system behaves as if the final configuration
had been assembled statically. \SZC{} gives this promise formal content: a
small-step calculus of components, fibers, and a guarded lifecycle, with a
metatheory proving preservation, recovery exactness, ordering, coherence,
progress, and confluence. Its \S 5 asserts that the Cordis runtime implements the
calculus, through a correspondence table (Table~2) and ten reference algorithms.
The assertion is not a theorem. Nothing in \SZC{} excludes an implementation that
matches Table~2 on names and violates Table~1 on behavior---and
\S\ref{sec:evaluation} exhibits a shipping build that does exactly that.

\subparagraph{The question.} What would it mean, mathematically, to \emph{verify}
that an implementation follows the calculus---and can a \emph{finite} set of
executable test cases carry that verification as a sufficient \emph{and} necessary
condition, rather than as anecdote?

\subparagraph{The obstruction.} For unrestricted implementations the answer is
no, by a classical diagonalization we make precise in
Theorem~\ref{thm:impossibility}: a black-box implementation may agree with the
specification on every name, key, and configuration size occurring in a finite
suite and deviate on a fresh one. Any honest finite characterization must
therefore name the structural hypothesis that closes this gap. The testing
literature has known this since Gaudel~\cite{gaudel95}: finite completeness is
always relative to \emph{uniformity} and \emph{regularity} hypotheses. Our
contribution is to give those hypotheses exact, auditable content for the
dynamic-composition calculus, and then to prove the biconditional they enable.

\subparagraph{The idea.} Each rule of the calculus reads its state through a
bounded set of \emph{logical atoms}: lifecycle-state tests, retirement, equality
of committed and target views, per-key satisfaction, existential witnesses for
conflict, reliance, and childlessness. Operations on unbounded structures
decompose pointwise (per key) or through witnesses (per edge). After names, keys,
values, and tags are normalized, each rule's window collapses to a finite
multiset of \emph{occurrences} drawn from a finite alphabet $\occ r$. A
\emph{clause-factored} implementation is one whose rule decisions are a lift of a
per-occurrence decision table $f_r^{\Impl} : \occ r \to \mathrm{Decision}_r$,
equivariantly and without hidden dependence on registry size, spelling, or
allocation order. For such implementations, one canonical experiment per
reachable signature suffices---and is necessary---to pin the entire behavior.

\subparagraph{Contributions.}
\begin{enumerate}
\item A formalization of the plugin-system interface of \SZC{}~\S5 as an
  instrumented LTS with events, refusals, transition certificates, quiescent
  reports, and divergence; conformance as termination-sensitive observation
  equivalence with rule-level bisimulation (\S\ref{sec:interface}).
\item An impossibility theorem for unrestricted finite testing, and the
  clause-factored implementation class \CF{} whose obligations CF1--CF7 are
  stated as checkable semantic properties, not slogans
  (\S\ref{sec:cf}).
\item The atomic-signature decomposition $\kap r : W_r \to \Mfin(\occ r)$ with
  fully enumerated alphabets---92 occurrences across the ten rules and the
  coeffect layer, 86 reachable with canonical experiments, 6 proved
  interface-unrealizable---and the finite characteristic suite \Tstar{}
  (\S\ref{sec:signatures}).
\item The characterization theorem
  $\Impl \in \CF \Rightarrow (\Impl \conf \Spec \iff \Impl \models \Tstar)$,
  proved by hand: factorization ($\Spec \in \CF$), transport, observation
  reflection, shortest-counterexample safety, and ranking-based termination
  (\S\ref{sec:theorem}). Two byproducts are of independent interest: a
  \emph{masking} proposition showing reflection is false in the
  certificate-free vocabulary (\S\ref{sec:reflection}), and a progress theorem
  that removes the acyclicity hypothesis of \SZC{} Thm.~66
  (\S\ref{sec:termination}).
\item An executable grounding: a reference semantics with an enumerable
  scheduling oracle (every test scenario's complete schedule tree explored; the
  largest has 924 schedules), a 27-member semantic deviant family with its kill
  matrix, and a bounded state-space explorer and Lean~4 mechanization in
  progress, reported with exact status (\S\ref{sec:mech}).
\item An evaluation of the current Cordis TypeScript implementation: one build
  passes all 32 obligations; the published \texttt{@deepseek-ai/cordis} 4.0.1
  fails exactly four, one previously unknown; a two-node distributed deployment
  fails exactly the same four (\S\ref{sec:evaluation}). Three corrections to
  \SZC{} itself (\S\ref{sec:discussion}).
\end{enumerate}

\subparagraph{Non-claims.} The proofs are rigorous mathematics in the style of
the source paper, not machine-checked derivations; the mechanization track and
its exact status are reported in \S\ref{sec:mech}. Membership of a concrete
codebase in \CF{} is an audit obligation, not a testable property; the
evaluation separates what the suite establishes from what the audit does.

%% =======================================================================
\section{Background and Motivation}\label{sec:background}

\subsection{The calculus of dynamic composition}

We recall from \SZC{} only what the development manipulates, in its notation.
A \emph{component} is a triple $(d, p, e)$: a coeffect specification
$d \subseteq K$ (keys read), a provision declaration $p \subseteq K$ (keys
written), and a witnessed effect function $e$ carrying its own inverse
(\SZC{} Defs.~8, 43). A \emph{fiber}
$\langle d, p, e, \pi, \sigma, \tau, \theta\rangle$ instantiates a component
under parent $\pi$ with its own table $\sigma$, retirement flag $\tau$, and
lifecycle state
\[
\theta \in \{\Inactive(\zeta),\ \Reloading(i,g,\omega),\ \Active(g,\omega),\
\Unloading(g,\omega,\zeta)\},
\]
where $g$ is the inverse accumulator, $\omega : d \to \mathfrak N$ the committed
resolution, $i$ the remaining effect iterator, and
$\zeta \in \{\bot\} \cup \Xi$ an outcome (\SZC{} Defs.~44, 49, 51). A state
$\gamma$ carries a registry $F_\gamma$ of named fibers; the derived coeffect
context is the union of \emph{Active} tables only,
$\sigma_\gamma = \bigcup \{\sigma_m \mid \theta_m = \Active(-,-)\}$, and the
target view $\target_n(\gamma)$ maps each declared key to its provider, or is
$\bot$ if $n$ is retired or a key lacks an Active provider (\SZC{} Def.~46).
Ten rules generate the semantics (\SZC{} Table~1): O\mbox{-}Insert (freshness;
live parent; declared disjointness $\forall m.\ p \cap p_m = \varnothing$),
O\mbox{-}Retire, O\mbox{-}Remove (retired, Inactive, childless), and seven
lifecycle rules: L\mbox{-}Begin commits $\omega$ where the target is defined;
L\mbox{-}Iter runs one iteration against the committed view, composing its
inverse LIFO; L\mbox{-}Finish reaches Active; L\mbox{-}Divert aborts a
transition whose target turned; L\mbox{-}Raise routes a failing iteration
through recovery to $\Inactive(\xi)$; L\mbox{-}Leave stops providing;
L\mbox{-}Unload---guarded by $\neg\relied_n(\gamma)$: no installed fiber still
resolves a key to $n$---applies the accumulator and discards $\omega$. An
iteration may \emph{register} a child, whose retirement is that iteration's
inverse (\SZC{} Def.~47). The metatheory we build on: rule invariance under the
observational equivalences (\SZC{} Lemma~55), equivariance (Lemma~56), framing
(Lemmas~54, 57), ordering and value fixity (Thm.~63), progress with the
measure $S(n) \le (K{+}4)(V(n){+}1)$ (Thm.~66), and confluence (Thm.~73).

\subsection{The verification gap}\label{sec:gap}

\SZC{}~\S5 maps the calculus onto Cordis: \texttt{ctx.plugin} for insertion,
\texttt{fiber.dispose} for retirement, \texttt{ctx.provide}/\texttt{get}/%
\texttt{set} for the coeffect operations, \texttt{ctx.effect} for tracked
effects, a committed per-fiber store with a proxy access walk (Algorithm~6) for
the read discipline, and a dependent-await (Algorithm~5) for the unload guard.
A prior targeted audit cataloged nine places where the shipped runtime and the
calculus disagree, five of which were fixed in a vendored build. What was
missing is the other direction: a principled argument that any finite battery of
checks \emph{could} certify conformance---and a battery provably sufficient and
necessary. That is what this paper constructs. Motivation is not hypothetical:
our evaluation (\S\ref{sec:evaluation}) shows the published build violating the
unload guard's ordering in precisely the way a deviant semantics predicts, and
leaking a recorded failure outcome across the fiber boundary on a
reconfiguration path no targeted audit had visited.

%% =======================================================================
\section{The Plugin-System Interface}\label{sec:interface}

\begin{definition}[events and observations]\label{def:events}
Fix countable disjoint sorts: names $\mathfrak N$, keys $K$, values
$\mathbb V$, effect tags $\mathbb T$. The event alphabet is
\[
\mathcal E ::= \mathsf{app}(n,t) \mid \mathsf{inv}(n,t) \mid \mathsf{rd}(n,k,v)
\mid \mathsf{rderr}(n,k,\epsilon) \mid \mathsf{act}(n) \mid \mathsf{deact}(n)
\mid \mathsf{val}(n,k)
\]
with $\epsilon \in \{\mathsf{IA}, \mathsf{UD}\}$ (inactive resp.\ undeclared
access). In addition, transitions emit \emph{certificates}
$\mathsf{cert}(r, \widehat s, \widehat a)$: the rule tag, the normalized
signature of the window it fired at, and the normalized update abstract. An
\emph{outcome} of a run is $\mathsf{Done}(w,q)$, $\mathsf{Refused}(w,\lambda,q)$,
or $\mathsf{Diverge}(u,v)$, where $w$ is the word of events and certificates,
$q$ the quiescent report (statuses, outcomes, the store
$\mathrm{st}(k) = \sigma_\gamma(k)$, refusal log), and $\lambda$ the refused
input. $\Out_{\mathbb X}(\rho)$ is the outcome set of system $\mathbb X$ on
orchestration script $\rho$ over all fair schedulers; $\equiv$ compares outcome
sets up to sort-preserving renaming and the calculus's value equivalence.
\end{definition}

\begin{definition}[plugin system; conformance]\label{def:conf}
A \emph{plugin system} is an LTS over orchestration labels
$\{\mathsf{insert}, \mathsf{retire}, \mathsf{update}, \mathsf{setval},
\mathsf{isolate}, \mathsf{intercept}\}$ and internal lifecycle steps, with the
instrumentation of Definition~\ref{def:events}. The specification \Spec{} is
the calculus so instrumented, with components given as \emph{scripts}: finite
step lists over
$\{\mathsf{provide}, \mathsf{track}, \mathsf{read}, \mathsf{setval},
\mathsf{register}, \mathsf{raise}, \mathsf{raiseUnless}\}$, one step per
iteration of \SZC{} Def.~51 (a tracked inverse may itself read---the
read-on-revert form---which \S\ref{sec:signatures} shows the suite requires).
$\Impl \conf \Spec$ iff for every script $\rho$:
(i) $\Out_\Impl(\rho) \equiv \Out_\Spec(\rho)$ restricted to
$\mathsf{Done}/\mathsf{Refused}$ outcomes, and
(ii) $\Impl$ has a $\mathsf{Diverge}$ outcome on $\rho$ only where \Spec{}
does. A test is $t = (\rho_t, V_t)$ with $V_t$ decidable on outcomes;
$\Impl \models t$ iff \emph{no run of $\Impl$ on $\rho_t$ diverges at a settle
point} and every outcome satisfies $V_t$.
\end{definition}

The divergence clause in $\models$ closes a vacuity: without it an
implementation that never quiesces passes every suite. Conformance is a
preorder; its bisimulation presentation (each tagged implementation transition
matched by a specification step with the same certificate and emission, and
conversely each specification-enabled step offered) is what the sufficiency
proof establishes and what necessity consumes.

%% =======================================================================
\section{Impossibility, and the Clause-Factored Class}\label{sec:cf}

\begin{theorem}[no finite suite characterizes unrestricted conformance]
\label{thm:impossibility}
For every finite suite $T$ there is an implementation $\Impl_T$ with
$\Impl_T \models T$ and $\Impl_T \not\conf \Spec$.
\end{theorem}

\begin{proof}[Proof (sketch; full proof in the artifact)]
Let $N$ bound the number of dependency keys any component declared in $T$'s
scripts. Let $\Impl_T$ behave exactly as \Spec{} except that L\mbox{-}Begin
never fires for fibers with $|d_n| > N$. Every observable of every run of every
$\rho_t$ is unchanged, so $\Impl_T \models T$; a script inserting a component
with $N{+}1$ satisfiable keys distinguishes it from \Spec. The same
construction works for fresh key spellings or registry sizes, so no finite
suite escapes it.
\end{proof}

The theorem locates exactly what a finite characterization must assume: that
the implementation cannot \emph{count}, \emph{spell}, or otherwise read its
state except through the rule's own atoms. We make this a definition.

\begin{definition}[clause-factored implementations]\label{def:cf}
$\Impl \in \CF$ iff:
\begin{description}
\item[CF1 (abstraction \& tagging)] there is a total abstraction
  $A : \Reach(\Impl) \to \Gamma_\Spec$ with $A(\gamma^0_\Impl) = \gamma^0$, and
  every implementation transition is tagged with the rule it implements;
  certificates are emitted per Definition~\ref{def:events}.
\item[CF2 (equivariance)] behavior commutes with sort-preserving bijections of
  names, keys, values, and tags fixing the reserved constants.
\item[CF3 (signature factorization)] for each rule $r$ there is a table
  $f_r^{\Impl} : \occ r \to \mathrm{Decision}_r$ such that tagged applicability,
  updates, and emissions at any window are the \emph{lift} of $f_r^{\Impl}$
  over the window's occurrence multiset $\kap r(W)$, combined by the row's
  declared folds (pointwise products over keys; existential folds over
  witnesses; one fixed binary composition for the accumulator), with order data
  drawn from the window's frame.
\item[CF4 (instrumentation faithfulness)] certificates report the actual
  $(\!r, \widehat s, \widehat a)$ of each tagged transition, and the
  certificate abstract $\widehat a$ covers the update aspects that events and
  reports do not determine (\emph{certificate adequacy};
  \S\ref{sec:reflection} shows this clause is not optional).
\item[CF5 (canonical reachability)] renaming-invariant behavior on the
  canonical experiments of \S\ref{sec:signatures}.
\item[CF6 (fair, exact settling)] the scheduler is fair and settle terminates
  exactly on quiescent configurations.
\item[CF7 (refusal discipline)] orchestration inputs are refused exactly where
  the rule's premises fail, with the refusal observable.
\end{description}
\end{definition}

Each obligation is auditable against code: CF3 rules out branching on
$|d_n|$, on registry cardinality, or on spellings (the Theorem~%
\ref{thm:impossibility} implementation violates it); CF2 rules out reserved
names; CF6 pins the scheduler. The class is far from vacuous: it contains
\Spec{} itself (Theorem~\ref{thm:fact}) and all twenty-seven deviant calculi of
\S\ref{sec:evaluation}, which is what makes necessity meaningful.

%% =======================================================================
\section{Atomic Signatures and the Characteristic Suite}\label{sec:signatures}

\subsection{Occurrence decomposition}

For each rule $r$, the \emph{window} $W_r(\gamma, n)$ is the tuple of fields
the Table-1 row reads and writes, together with its \emph{combination frame}:
the index sets of its finite-map operations and the order data of its
order-sensitive operators (iterator position; accumulator order). The
normalization
\[
\kap r : W_r \to \Mfin(\occ r)
\]
maps a window to a finite multiset of \emph{occurrences}: one scalar occurrence
(the row's scalar guard-atom valuation, step constructor, and equality
pattern); one per-key occurrence per key of each pointwise index set; one
witness occurrence per satisfied existential (conflict key, reliance edge,
child edge); one per-item occurrence per accumulator entry. Crucially,
$\kap r$ keeps equality \emph{patterns} and forgets identities and counts:
a registry of a thousand independent fibers contributes nothing to a rule that
does not read them (framing), and a dependency set of any size contributes only
the set of satisfaction \emph{classes} present.

\subsection{The enumerated alphabets}

The complete tables---every occurrence with its atoms, a finiteness proof per
alphabet, reachability analysis, and a canonical experiment per reachable
signature---occupy the artifact's \texttt{signatures.md} (1{,}342 lines);
Table~\ref{tab:sig} summarizes.

\begin{table}[t]
\caption{Occurrence alphabets: sizes, reachable subsets, canonical scripts.}
\label{tab:sig}
\centering
\small
\begin{tabular}{lccc}
\toprule
rule / operation & $|\occ r|$ & reachable & scripts \\
\midrule
O-Insert & 7 & 6 & 4 \\
O-Retire & 4 & 4 & 4 \\
O-Remove & 7 & 7 & 5 \\
L-Begin & 7 & 7 & 6 \\
L-Iter & 11 & 11 & 8 \\
L-Finish & 3 & 3 & 3 \\
L-Divert & 6 & 4 & 2 \\
L-Raise & 9 & 9 & 7 \\
L-Leave & 5 & 4 & 3 \\
L-Unload & 12 & 12 & 9 \\
coeffect ops (Defs.~23/24, 28/29, 31; Alg.~6) & 21 & 19 & 9 \\
\midrule
total & 92 & 86 & 60 \\
\bottomrule
\end{tabular}
\end{table}

Six occurrences are proved \emph{interface-unrealizable}, each with an argument
rather than an assertion: O\mbox{-}Insert freshness violation (names are
system-drawn); the L\mbox{-}Divert landing alternative (requires the
asynchrony layer of \SZC{}~\S4.3.3, absent from the synchronous reference
semantics); the changed-provider divergence class at both L\mbox{-}Divert and
L\mbox{-}Leave---ruled out by a lemma: a live committed edge pins its provider,
by the reliance guard together with declared disjointness, so a target can
change only \emph{to or from} $\bot$; and two Algorithm-6 access classes
(self-inactive; withdrawn-binding) blocked by view totality and the guard,
respectively---the latter being exactly the content of \SZC{} Thm.~63
reappearing as an unreachability fact.

\subsection{Canonical experiments and \Tstar}

For each rule $r$ and reachable signature $s$, the canonical experiment
$(\rho_{r,s}, V_{r,s})$ drives the system, by an explicit orchestration script
over explicit component data, to a pointed configuration realizing $s$, lets
the attempt at $r$ occur (directly for orchestration rules; by settling for
lifecycle rules), and asserts $\equiv$-equality of the implementation's outcome
set with the specification's at that point---certificates included. Then
\[
\Tstar = \{(\rho_{r,s}, V_{r,s}) \mid r \in \mathrm{Rules} \cup \mathrm{Ops},\
s \in \occ r^{\mathrm{reach}}\}, \qquad |\Tstar| = 60 \text{ scripts}.
\]
Some scripts are load-bearing in nontrivial ways: realizing the
\emph{installed-only-provided} satisfaction class requires observing a begin
attempt inside another fiber's guarded teardown window; realizing the
inactive-access error under Algorithm~6 requires five components and five
settles. The 32-test obligation suite that predates this construction
(\S\ref{sec:evaluation}) is \emph{material} for \Tstar: the explorer track
(\S\ref{sec:mech}) audits which signatures it covers, and the audit---not
intuition---decides whether it must grow.

%% =======================================================================
\section{The Characterization Theorem}\label{sec:theorem}

\begin{theorem}[finite characterization]\label{thm:main}
For every $\Impl \in \CF$:
\(
\Impl \conf \Spec \iff \Impl \models \Tstar,
\)
where observation equivalence is taken in the certificate vocabulary of
Definition~\ref{def:events}.
\end{theorem}

The proof decomposes into five results, each proved in full in the artifact's
\texttt{proofs.md} (546 lines); we state them and give the proof shapes,
displaying the two arguments of independent interest.

\subsection{Factorization: the calculus is clause-factored}

\begin{theorem}[$\Spec \in \CF$]\label{thm:fact}
Each rule's applicability is a function of its window's occurrence multiset;
each write is the lift of per-occurrence updates under the row's declared
folds; emissions are per-occurrence; equivariance and framing hold.
\end{theorem}

\begin{proof}[Proof shape]
Rule by rule against Table~1: the guard atoms are read through the occurrence
scalar; target construction is a pointwise product of singleton-key provider
decisions; conflict, reliance, and childlessness are existential folds over
witness occurrences; the accumulator is a fold of one binary operator whose
order lives in the frame. Equivariance extends \SZC{} Lemma~56 from names to
all four sorts by a per-premise inventory; framing packages Lemmas~54(1)
and~57. A symbol inventory per row verifies that \emph{no rule reads a
cardinality}: Table~1 contains no arithmetic. This last observation is the
exact property Theorem~\ref{thm:impossibility}'s counterexample lacks.
\end{proof}

\subsection{Transport}

\begin{lemma}[transport]\label{lem:transport}
For $\Impl \in \CF$: windows with equal occurrence multisets (up to
sort-preserving renaming) receive equal per-occurrence decisions; whole-rule
results additionally correspond when the windows' \emph{frames} correspond.
\end{lemma}

The second clause is not pedantry: accumulator \emph{order} lives in the frame,
and a whole-decision transport across order-differing, multiset-equal windows
is false. (This corrected an earlier draft of the theory in which order data
was homeless.)

\subsection{Reflection, and the necessity of certificates}\label{sec:reflection}

\begin{lemma}[observation reflection]\label{lem:reflection}
If $f_r^{\Impl}(s) \neq f_r^{\Spec}(s)$ at a reachable signature $s$, then
$\Impl$ fails $(\rho_{r,s}, V_{r,s})$.
\end{lemma}

\begin{proof}[Proof shape]
By cases on the decision variants (spurious fire; missed fire; wrong update;
wrong emission; wrong refusal). The key device: at the realized signature, the
specification's certificate is forced to $f_r^{\Spec}(s)$ because the table is
a function---so a deviant certificate matches nothing in the specification's
outcome set, and $V_{r,s}$'s $\equiv$-equality fails. CF4 is used exactly where
the deviation would otherwise be invisible.
\end{proof}

\begin{proposition}[masking]\label{prop:masking}
In the certificate-free vocabulary (events and reports only), reflection is
false: there is a rule aspect---skipping provision inverses at a
\emph{retirement} unload---whose deviation is unobservable on \emph{every}
script, because removal deletes the entry any residue would inhabit; moreover
the deviant \emph{conforms} in that vocabulary.
\end{proposition}

The proposition rescues the theorem honestly rather than weakening it silently:
the biconditional holds in the certificate vocabulary, and holds in the
certificate-free vocabulary only for the quotient of \Spec{} by the masked
aspects. It also explains an empirical fact we had met earlier: mutation
analysis of the concrete suite found exactly this aspect unkillable.

\subsection{Safety}\label{sec:safety}

\begin{theorem}[safety direction]\label{thm:safety}
If $\Impl \in \CF$ and $\Impl \models \Tstar$, every finite behavior of
$\Impl$ is a behavior of \Spec.
\end{theorem}

\begin{proof}[Proof shape]
Minimal counterexample over the first deviating index of a run (both
deviation polarities: a step \Spec{} cannot match; a specification-enabled step
$\Impl$ cannot offer). The matched prefix gives the reachability bridge: the
deviating window's signature is reachable \emph{in \Spec} because the prefix
agrees; Lemma~\ref{lem:transport} transports the deviation to the canonical
realization $\rho_{r,s}$; Lemma~\ref{lem:reflection} makes it fail there,
contradicting $\Impl \models \Tstar$.
\end{proof}

\subsection{Termination, and progress without acyclicity}\label{sec:termination}

\begin{theorem}[generalized progress; strengthens \SZC{} Thm.~66]
\label{thm:progress}
Every maximal lifecycle sequence of \Spec{} is finite---\emph{without} the
acyclicity hypothesis on the provider precedence relation $\prec$.
\end{theorem}

\begin{proof}[Proof shape]
An $\omega$-provenance invariant shows every committed edge points at a fiber
installed at commitment time; a \emph{frozen-cycle} lemma then shows no member
of a co-present $\prec$-cycle is ever installed (first-co-presence freshness
plus instantaneous propagation of well-formedness, \SZC{} Def.~58(2,4)); a
frozen cone extends this to everything a cycle supports. Outside the frozen
cone, the per-fiber and attribution counting bounds internal to \SZC{}
Thm.~66's proof are re-proved standalone, and infinite descent closes the
argument. The calculus thus quiesces on cyclic configurations not by the
theorem's bound but by unsatisfiability---the cycle freezes.
\end{proof}

\begin{theorem}[termination direction]\label{thm:term}
For $\Impl \in \CF$ with $\Impl \models \Tstar$: $\Impl$ diverges only where
\Spec{} does. Untagged stuttering is excluded by CF1's tagging totality with
CF6; the ranking of Theorem~\ref{thm:progress} transfers along the step
matching of Theorem~\ref{thm:safety}.
\end{theorem}

Termination-sensitivity is not decorative: necessity survives its deletion,
sufficiency does not (two counterexamples in the artifact).

\subsection{Assembly}

Necessity: if $\Impl \conf \Spec$, each canonical experiment's outcome set
equals \Spec's by the bisimulation, so every $V_{r,s}$ passes and no settle
diverges. Sufficiency: Theorems~\ref{thm:safety} and~\ref{thm:term} assemble
the observation bisimulation clause by clause. \hfill$\qed$

%% =======================================================================
\section{Mechanization and Exhaustive Exploration}\label{sec:mech}

Three executable artifacts ground the theory; we report exact status.

\subparagraph{Exhaustive schedules (complete).} The reference semantics
implements the calculus's nondeterminism as a \emph{choice oracle}: wherever
more than one rule is applicable, the oracle picks. A mixed-radix enumerating
driver explores every oracle sequence of every concrete test scenario. Of the
32 obligation scenarios, 25 are schedule-deterministic; the largest choice tree
has 924 schedules; all trees are enumerated to completion. Soundness of the
suite on \Spec{} and the entire deviant matrix of \S\ref{sec:evaluation} are
certified under this enumeration---the schedule quantifier in ``$\models$'' is
discharged exactly, not sampled.

\subparagraph{Bounded state-space explorer (in progress).} A normalized
all-successor BFS over the reference semantics at bound (4 fibers, 2 keys,
5 steps) with canonicalization of names, keys, and values, generating: the
reachable signature sets per rule (an independent check of
Table~\ref{tab:sig}'s reachability column, including its six negative claims),
a generated bounded suite $T_B$ with shortest access sequences, the coverage
matrix of the existing 32 tests against reachable signatures, and the deviant
kill matrix of the generated suite. At the time of writing the pipeline
(normalization with self-test, audit, generation, kill-matrix) is validated
and the full-catalog exploration runs are executing; results will accompany
the artifact.

\subparagraph{Lean~4 mechanization (in progress).} The bounded theory in core
Lean (no external libraries): the calculus at a finite bound as a decidable
successor function; equivariance; finiteness of the configuration type;
explorer soundness and completeness; the bounded characterization by
reflection; the transport lemma as the \emph{single named axiom} whose paper
proof is Lemma~\ref{lem:transport}; and the conditional final theorem. Status
at submission is recorded in the artifact's \texttt{PROOF\_STATUS.md}; nothing
in \S\ref{sec:theorem} depends on it---the dependency runs the other way.

%% =======================================================================
\section{Evaluation of the Cordis TypeScript Implementation}
\label{sec:evaluation}

\subsection{Targets and instruments}

We evaluate with the concrete 32-obligation suite (the predecessor and
material of \Tstar), each obligation pinned to a clause of the calculus, under
the exhaustive-schedule harness, against:
(i) the reference semantics (soundness control);
(ii) \emph{Cordis-aligned}: the vendored build carrying the prior audit's
calculus-alignment fixes;
(iii) \emph{Cordis-upstream}: \texttt{@deepseek-ai/cordis}~4.0.1 as published;
(iv) the upstream build deployed as \emph{two nodes} joined by the
\texttt{cordis-node} projection layer (in-process transports, mutual mounts,
component placement alternating across nodes, observation through projection).
Necessity is instrumented by a family of 27 \emph{semantic deviants} of the
reference semantics, each negating one clause (guard dropped; accumulator
FIFO; $\sigma_\gamma$ over installed fibers; re-entry from $\Inactive(\xi)$;
live-view reads; \dots)---all clause-factored, hence inside the class the
theorem quantifies over.

\subsection{Results}

\begin{table}[t]
\caption{Suite results (32 obligations; schedule trees fully enumerated).}
\label{tab:results}
\centering
\small
\begin{tabular}{lc}
\toprule
target & result \\
\midrule
reference semantics & 32/32 \\
Cordis-aligned build & \textbf{32/32} \\
Cordis-upstream 4.0.1 & 28/32 \\
two-node distributed (upstream) & 28/32 --- the same four \\
\midrule
deviant family (27) & all killed; 14 obligations essential \\
\bottomrule
\end{tabular}
\end{table}

\subparagraph{The four upstream divergences.}
\emph{(1, 2) Unload-guard ordering} (two obligations): provider inverses run
concurrently with, and before, dependents' guarded teardown---the guard is
implemented one level too low (inside the provision effect's own disposer,
while the fiber's unload starts all disposers at once). The observed event
words realize the guard-deviant derivation of the necessity analysis nearly
letter for letter. \emph{(3) Declared disjointness} unenforced at insertion.
\emph{(4) Escaped failure outcome} (previously unknown): after a failed first
activation, a later \texttt{update()} of the by-then recovered, Active fiber
releases the \emph{first} episode's error as a process-level unhandled
rejection; the calculus requires $\xi$ recorded on the fiber and nothing to
escape. States and effect traces are otherwise correct, which is why no
state-based audit had found it. All four are fixed in the aligned build, whose
32/32 stands under complete schedule enumeration.

\subparagraph{What 32/32 does and does not claim.} By
Theorem~\ref{thm:main}, passing certifies conformance \emph{given}
$\Impl \in \CF$; the CF obligations are audit items on the TypeScript source
(no cardinality branching, per-key folds, one accumulator operator, fair
scheduler), constructed for the abstraction clause and asserted---not
proved---for the rest. The evaluation therefore reads: \emph{the aligned build
is conformant relative to a stated, checkable audit obligation; the published
build is non-conformant outright} (a single failed obligation refutes
membership of the conformant class with no further hypothesis).

\subparagraph{Distribution.} Under the equivalence contract of the
\texttt{cordis-node} projection layer (established in its own companion
verification), conformance transfers to distributed deployments by
composition. Empirically, the two-node deployment fails exactly the upstream
build's four obligations: the difference attributable to distribution on this
suite is empty---with genuine cross-node execution (provider on one node,
dependent on the other, satisfied through projection).

%% =======================================================================
\section{Discussion}\label{sec:discussion}

\subparagraph{Findings about the source paper.} The construction returned
three corrections to \SZC{}: \emph{(PB-1)} Table~2 maps Def.~23's
$\mathrm{set}(k,v)$ to \texttt{ctx.set}; the primitive carrying that
operation's precondition and inverse is \texttt{ctx.provide}, while
\texttt{ctx.set} is a value operation on an existing binding.
\emph{(PB-2)} for in-place overwrite by an Active provider, \S5.1.3 prescribes
``not observed'' while Thm.~63(3) forces observed replacement if the operation
is admitted at all; the calculus is silent, the two prose commitments
contradict, and our suite pins the dichotomy any extension must satisfy (never
a torn episode). \emph{(PB-3)} Algorithm~10 replaces a component by
\texttt{dispose()} immediately followed by \texttt{use()}, racing
O\mbox{-}Insert's own disjointness premise; the reference semantics loses that
race under adversarial scheduling, latching a spurious failure---the staged
retire--drain--insert is the form whose premises hold at every step.

\subparagraph{What the theory added beyond testing folklore.} Three things.
The impossibility theorem plus \CF{} converts ``we tested it'' into a
two-part claim with an explicit audit residue. The masking proposition shows
that observation vocabularies are part of the theorem statement: without
certificates, some deviations are not just unkilled but \emph{unkillable}, and
the correct response is a vocabulary extension, not a bigger suite. And the
frozen-cycle progress theorem shows the calculus is better-behaved than its
own metatheory claimed: cyclic dependency configurations quiesce by freezing,
so the acyclicity hypothesis of \SZC{} Thm.~66 can be deleted.

\subparagraph{Limitations.} Membership in \CF{} is assumed of implementations
and audited, not tested---Theorem~\ref{thm:impossibility} shows this is not a
weakness of our suite but of finite testing. The hand proofs are unmechanized;
the Lean track isolates precisely one axiom (transport) and is in progress.
The six unrealizability proofs and the reachability column await independent
confirmation by the explorer. The evaluation's schedule exhaustion applies to
the reference semantics; for the TypeScript builds, schedules are exercised
through the real event loop and quantified by CF6. The distributed result is
conditional on the projection layer's own contract.

%% =======================================================================
\section{Related Work}\label{sec:related}

\subparagraph{Testing from formal specifications.} Gaudel's
theory~\cite{gaudel95} established finite completeness relative to uniformity
and regularity hypotheses; \CF{} is those hypotheses made concrete for a
lifecycle calculus, and Theorem~\ref{thm:impossibility} is the classical
diagonalization in our setting. Tretmans' \textsf{ioco}
theory~\cite{tretmans96, tretmans08} gives sound and exhaustive suites for
LTSs with quiescence; our settle observation plays $\delta$'s role, and our
proofs take a simulation form because the calculus's own invariance lemma is
natively an up-to technique. Complete test suites for FSM equivalence (the
W-method~\cite{chow78}) inspire the characterization shape; our setting
differs by unbounded state (hence signatures rather than state counting),
nondeterminism (hence outcome sets over fair schedulers), and refusals.

\subparagraph{Mutation analysis.} Adequacy against deviants follows
DeMillo--Lipton--Sayward~\cite{demillo78}; mutating an executable
\emph{specification} follows specification-mutation work~\cite{black00};
subsumption among tests is documented in the surveys~\cite{jia11,
papadakis17}---our masking proposition gives one mechanism of unkillability a
proof.

\subparagraph{Executable semantics.} Validating a semantics by executing it
against suites is the semantics-engineering methodology~\cite{felleisen09}
and the K framework's practice~\cite{rosu10}; our reference interpreter with
an enumerable scheduling oracle is that methodology plus exhaustive schedule
coverage.

\subparagraph{Conformance of module and plugin systems.} De-facto conformance
suites (e.g.\ ECMAScript's test262) carry no completeness theorem; mechanized
module-system metatheory verifies calculi, not implementations' fidelity to
them. We are not aware of a prior finite characterization theorem for a
dynamic-composition runtime.

%% =======================================================================
\section{Conclusion}\label{sec:conclusion}

The question ``does this implementation follow the calculus?'' has an exact
answer. Unrestricted, no finite suite answers it. Restricted to
clause-factored implementations---a class whose obligations are auditable
against source code and whose definition is forced by the impossibility
theorem---a suite of one canonical experiment per reachable atomic signature
is both sufficient and necessary:
$\Impl \conf \Spec \iff \Impl \models \Tstar$. The signature alphabets are
finite and enumerated (92 occurrences; 86 reachable; 6 proved unrealizable);
the proof is by hand in the source paper's own style, and en route strengthens
that paper's progress theorem and corrects three of its claims. Applied to
Cordis, the theory locates four defects in the published build---one invisible
to state-based auditing---certifies the aligned build relative to a stated
audit residue, and shows the distribution layer preserving the verdict
wholesale. The remaining distance to a machine-checked artifact is a single
named axiom and a bounded model, both under construction; the mathematics no
longer waits on them.

%% =======================================================================
\appendix

\section{Artifact and Reproduction}\label{app:artifact}

The artifact accompanies the paper in the same repository:
\texttt{formal/signatures.md} (the 92-occurrence tables, finiteness and
reachability proofs, 60 canonical scripts);
\texttt{formal/proofs.md} (the complete proofs of
Theorems~\ref{thm:fact}--\ref{thm:term} and the assembly);
\texttt{formal/explorer/} (the bounded explorer) and
\texttt{formal/lean/} (the mechanization), each with status files;
\texttt{proof/} (the executable reference semantics, the 32-obligation suite,
the 27 deviants, and the harnesses). Reproduction:

\begin{verbatim}
cd paper/proof
MODEL_EXHAUST=1 node --test tests/*.test.mjs   # all schedules
PROOF_TARGET=cordis node --test tests/*.test.mjs
PROOF_TARGET=cordis CORDIS_LIB=<build>/lib/index.js \
  node --test tests/*.test.mjs
node run-necessity.mjs                          # deviant matrix
PROOF_TARGET=cordis-node node --test tests/*.test.mjs
\end{verbatim}

Node $\ge 22$; no dependencies; the deviant matrix's exit code certifies
adequacy and baseline soundness.

\section{Signature Table Excerpt}\label{app:sig}

By way of sample, the L\mbox{-}Unload alphabet (12 occurrences): the scalar
occurrence records $(\theta_n = \Unloading$, $\relied$-witness present?,
outcome class $\in \{\bot, \xi\}$, $\tau$, post-target class
$\in \{\bot, \text{satisfiable}\})$ with the well-formed $\tau \times$
post-target blocks reduced to three; witness occurrences record each reliance
edge (installed dependent, committed key, equality pattern); per-item
occurrences record each accumulator entry's inverse class
$\in \{\mathsf{provide}^{-1}, \mathsf{track}^{-1}, \mathsf{register}^{-1},
\mathsf{read\mbox{-}on\mbox{-}revert}\}$. The guard-blocked signature's
canonical experiment retires a provider under an installed consumer whose
tracked inverse reads the provided key; its verdict asserts the consumer's
read event, its precedence over every provider inverse, and the final
disposal---the experiment that convicts the published build.

\begin{thebibliography}{9}

\bibitem{gaudel95}
M.-C. Gaudel.
\newblock Testing can be formal, too.
\newblock In \emph{TAPSOFT '95}, LNCS 915, 1995.

\bibitem{tretmans96}
J. Tretmans.
\newblock Test generation with inputs, outputs and repetitive quiescence.
\newblock \emph{Software---Concepts and Tools} 17(3), 1996.

\bibitem{tretmans08}
J. Tretmans.
\newblock Model based testing with labelled transition systems.
\newblock In \emph{Formal Methods and Testing}, LNCS 4949, 2008.

\bibitem{chow78}
T. S. Chow.
\newblock Testing software design modeled by finite-state machines.
\newblock \emph{IEEE TSE} 4(3), 1978.

\bibitem{demillo78}
R. A. DeMillo, R. J. Lipton, F. G. Sayward.
\newblock Hints on test data selection: Help for the practicing programmer.
\newblock \emph{IEEE Computer} 11(4), 1978.

\bibitem{black00}
P. E. Black, V. Okun, Y. Yesha.
\newblock Mutation operators for specifications.
\newblock In \emph{ASE 2000}.

\bibitem{jia11}
Y. Jia, M. Harman.
\newblock An analysis and survey of the development of mutation testing.
\newblock \emph{IEEE TSE} 37(5), 2011.

\bibitem{papadakis17}
M. Papadakis, Y. Jia, M. Harman, et al.
\newblock Mutation testing advances: An analysis and survey.
\newblock \emph{Advances in Computers}, 2017.

\bibitem{felleisen09}
M. Felleisen, R. B. Findler, M. Flatt.
\newblock \emph{Semantics Engineering with PLT Redex}.
\newblock MIT Press, 2009.

\bibitem{rosu10}
G. Ro\c{s}u, T. F. \c{S}erb\u{a}nu\c{t}\u{a}.
\newblock An overview of the K semantic framework.
\newblock \emph{J. Logic and Algebraic Programming} 79(6), 2010.

\end{thebibliography}

\end{document}
